\documentclass[10pt]{article}

\usepackage[letterpaper,margin=0.62in,headheight=14pt]{geometry}
\usepackage{amsmath,amssymb,bm,mathtools}
\usepackage{booktabs,tabularx,array}
\usepackage{xcolor,microtype,enumitem}
\usepackage[hidelinks]{hyperref}

\definecolor{routeblue}{RGB}{24,67,112}
\definecolor{routegreen}{RGB}{37,104,71}
\definecolor{routegray}{RGB}{75,83,92}
\definecolor{routeorange}{RGB}{153,82,20}

\newcommand{\vect}[1]{\bm{#1}}
\newcommand{\dd}{\,\mathrm d}
\newcommand{\R}{\mathbb R}
\newcolumntype{L}[1]{>{\raggedright\arraybackslash}p{#1}}
\newcolumntype{Y}{>{\raggedright\arraybackslash}X}
\newcommand{\MethodTitle}[2]{%
  {\Large\bfseries\color{routeblue}#1}\hfill
  {\normalsize\bfseries\color{routegray}#2}\par\vspace{2pt}
  \hrule\vspace{5pt}}

% Replace these three fields.
\newcommand{\NoteTitle}{Short scientific methods briefing}
\newcommand{\NoteSubtitle}{Consistent equations, implementation requirements, and evidence gates}
\newcommand{\NoteStatus}{Discussion draft --- \today}

\setlength{\parindent}{0pt}
\setlength{\parskip}{2.5pt}
\setlength{\abovedisplayskip}{4pt}
\setlength{\belowdisplayskip}{4pt}
\setlength{\jot}{2pt}
\setlist[itemize]{leftmargin=1.18em,itemsep=1.5pt,topsep=2pt,parsep=0pt}
\setlist[enumerate]{leftmargin=1.45em,itemsep=1.5pt,topsep=2pt,parsep=0pt}
\renewcommand{\arraystretch}{1.10}
\raggedbottom
\pagestyle{plain}

\begin{document}
\small

% ---------- Overview page ----------
\begin{center}
  {\LARGE\bfseries\color{routeblue}\NoteTitle}\\[3pt]
  {\large\NoteSubtitle}\\[4pt]
  {\normalsize\NoteStatus}
\end{center}
\vspace{-3pt}

\textbf{Objective.} State the prediction, inference, or decision problem in two or three sentences. Name the reference truth and the intended deployment regime.

\textbf{Common forward model.} Define the state $u(\vect x,t)$, parameter vector $\theta$, flux $\vect F$, source $s$, domain $\Omega$, initial state $u_0$, and boundary operator $\mathcal B$:
\begin{equation}
 \partial_t u+\nabla\cdot\vect F(u,\nabla u;\theta)=s(u,\vect x,t;\theta),
 \quad u(\vect x,0)=u_0(\vect x),
 \quad \mathcal B[u]=g\ \text{on }\partial\Omega .
 \label{eq:reference-model}
\end{equation}
Explain whether Eq.~\eqref{eq:reference-model} is exact, a candidate reduction, or a closure.

\textbf{Consistent notation.}
\begin{center}
\begin{tabularx}{0.96\textwidth}{L{0.18\textwidth}Y}
\toprule
Symbol & Meaning and units \\
\midrule
$\vect x,t$ & Spatial position and physical time. \\
$u,u_0$ & State and initial state; distinguish intensive and extensive forms if relevant. \\
$\theta$ & Known or inferred model coefficients; identify which components are positive or spatially varying. \\
$\vect F,s$ & Conservative flux and volumetric/source contribution. \\
$\mathcal H[u]$ & Observation or output operator used for comparison with data. \\
\bottomrule
\end{tabularx}
\end{center}

\textbf{Method map.}
\begin{center}
\begin{tabularx}{0.98\textwidth}{>{\bfseries\raggedright\arraybackslash}p{0.18\textwidth}Y Y}
\toprule
Method & Scientific question & Primary output \\
\midrule
Method A & What is retained and what is reduced? & State, field, or quantity of interest. \\
Method B & Which correction, coefficient, or operator is inferred? & Prediction and uncertainty, if applicable. \\
\bottomrule
\end{tabularx}
\end{center}

\textbf{Decision order.} State which method is the first test, what evidence advances the project, and what failure or independent need justifies the next method.

% ---------- Duplicate this page once per method ----------
\clearpage
\MethodTitle{Method A: descriptive name}{Primary role}

\textbf{Purpose.} Explain the method's job and what information it intentionally retains or discards.

\textbf{Forward equation.} Write the complete forward problem, not only its differential operator:
\begin{equation}
 \boxed{\mathcal M\,\partial_t q+\mathcal A(q;\theta)=f(t;\theta),
 \qquad q(0)=q_0,\qquad \mathcal B_q[q]=g_q.}
 \label{eq:method-a}
\end{equation}
Define $q$, $\mathcal M$, $\mathcal A$, $f$, $q_0$, and the boundary data. Explain each non-obvious term and mark fitted closures explicitly.

\textbf{Inputs and outputs.} List known geometry, coefficients, forcing, initial state, boundary conditions, and requested outputs. If an observation map is used, write
\begin{equation}
 y=\mathcal H[q]+\varepsilon
\end{equation}
and define the error or discrepancy model rather than calling all mismatch noise.

\textbf{Implementation requirements.}
\begin{itemize}
 \item Identify the required data products, metadata, and numerical operators.
 \item State conservation, positivity, uniform-state, stability, or dimensional checks.
 \item Name the independent-case and full-rollout validation needed for acceptance.
\end{itemize}

\textbf{Merit and risk.} Give a realistic assessment of physical fidelity, implementation effort, CPU/GPU needs, likely accuracy regime, and principal failure mode. State the evidence threshold for adding complexity.

% ---------- A second method page shows the comparison pattern ----------
\clearpage
\MethodTitle{Method B: descriptive name}{Alternative or escalation}

\textbf{Purpose.} State why this method exists after Method A and which documented residual or unmet need it addresses.

\textbf{Forward or learned map.}
\begin{align}
 \dot{\vect a}&=\vect f_r(\vect a,\mu,t;\theta),
 &\vect a(0)&=\mathcal E[u_0],\label{eq:reduced-state}\\
 \widehat u(\vect x,t)&=\mathcal D[\vect a(t)] .
\end{align}
Define the encoder $\mathcal E$, decoder $\mathcal D$, reduced state $\vect a$, parameters $\mu$, and learned or projected coefficients $\theta$.

\textbf{Evidence and constraints.} Identify invariant constraints, training-case separation, forecast horizon, and out-of-distribution conditions. Separate one-step fit from autonomous rollout accuracy.

\textbf{Implementation requirements.}
\begin{itemize}
 \item Give the smallest credible baseline and the next escalation only if it fails.
 \item State offline data, memory, CPU/GPU, and online evaluation requirements.
 \item Define failure in observable scientific quantities, not only training loss.
\end{itemize}

\textbf{Merit and risk.} Explain when this method can outperform Method A, what guarantees are lost, and how confidence should be calibrated before validation.

% ---------- Optional final decision page ----------
\clearpage
\MethodTitle{Comparison and recommended sequence}{Decision page}

\begin{center}
\begin{tabularx}{\textwidth}{>{\bfseries\raggedright\arraybackslash}p{0.18\textwidth}L{0.14\textwidth}Y Y}
\toprule
Method & Online cost & Principal merit & Present limitation \\
\midrule
Method A & Low/moderate & Physics or interpretability advantage. & State the evidence gap. \\
Method B & Low after training & Flexibility or acceleration advantage. & State data and extrapolation risk. \\
\bottomrule
\end{tabularx}
\end{center}

\textbf{Recommended sequence.}
\begin{enumerate}
 \item Establish data alignment, units, initial/boundary conditions, and numerical truth.
 \item Implement the simplest conservative or otherwise invariant-preserving baseline.
 \item Add parameters, closure, or learned capacity only after a held-out failure is diagnosed.
 \item Validate complete cases and uninterrupted rollouts; report scientific errors and compute cost.
\end{enumerate}

\textbf{Evidence status.} End with a precise statement of what is demonstrated, what remains a hypothesis, and which data or experiment controls the next decision.

\end{document}
